\begin{claim}{B.33}{Proved}{Executed positive catalogues in genera two and three}
For $(g,\nu)=(2,0),(2,1),(2,2),(3,0)$, the numbers of conjugacy
classes of finite orientation-preserving subgroups, allowing conjugacy
in the full extended mapping class group, are respectively
\[
                         21,\qquad10,\qquad23,\qquad55.
\]
The trivial subgroup is included. The files
\texttt{oriented\_higher\_genus\_catalogue.csv} and
\texttt{oriented\_higher\_genus\_sections.json} give each class,
its actual finite permutation group, every subgroup's number of
sections, all nontrivial partial-section classes and their inclusion
relations, the $G$-set $\Phi$, its stabilizers, and every value of
$\mu$ on every section. Punctures are unlabelled. These assertions
concern positive subgroups; they do not enumerate subgroups containing
orientation-reversing elements.
\end{claim}
\begin{proof}
We specify the classification input and the computation using it.
Karab\'a\v{s}--Nedela--Skyvov\'a
\cite[Theorem 2.2, Corollary 2.3, Theorem 4.1 and \S8]{R-KNS}
identify the positive topological classes with the orbits of smooth
epimorphisms from the Fuchsian group
\[
 \Delta=\langle x_1,\ldots,x_r,a_1,b_1,\ldots,a_h,b_h\mid
 x_i^{m_i}=1,\ [a_1,b_1]\cdots[a_h,b_h]x_1\cdots x_r=1\rangle
\]
under target automorphisms and orientation-preserving orbifold
automorphisms. Here smooth means that every image of $x_i$ has
order exactly $m_i$. Their \S8 gives a complete genus-two list
of 21 classes and a genus-three list of 55 classes in 49
group/signature pairs, including the trivial actions. We opened
the accompanying data \cite{R-KNSData}, downloaded the two original
JSON files, and retained their attribution and SHA-256 digests.
The attached program reads those exact files. It verifies that
each generating vector generates the group of the stated order,
each elliptic generator has its specified order, every presentation
relation holds, and
\[
 |G|\left(2-2h-\sum_i(1-1/m_i)\right)=2-2g.
\]
The conversion from the source's right-acting permutations to the
left-acting permutations in our code is the isomorphism sending
each permutation to its inverse. For a handle word the two
commutator conventions are compared explicitly; replacing both
handle generators by their inverses converts $a^{-1}b^{-1}ab$
to the convention $aba^{-1}b^{-1}$. The files retain the actual
generating vector, so this convention does not affect a recorded
elliptic rotation.

The local-order checks imply that the kernel of each epimorphism
is torsion-free: every finite-order element of a Fuchsian group
is conjugate to a power of an elliptic generator. The covering
formula then identifies its quotient surface as genus $g$.
Conversely a positive finite action determines this epimorphism
by lifting paths in the quotient. The two constructions are
inverse modulo the stated source and target automorphisms.
Finally C.5--C.6 identify these action classes with conjugacy
classes of finite mapping-class subgroups. This supplies the
exact reduction from the classification input to the groups
asked for here.

We also recompute the equivalence relation whenever $h=0$,
rather than using only the published class count. For an ordered
generating vector $(u_1,\ldots,u_r)$ with product one, the braid
generator acts by
\[
 \sigma_i:(u_i,u_{i+1})\longmapsto
               (u_i u_{i+1}u_i^{-1},u_i).
\]
Use adjacent such generators when $m_i=m_{i+1}$ and all pure
braid generators, obtained by moving the $j$th entry next to the
$i$th, applying $\sigma_i^2$, and reversing the moves.
Theorem 4.1 of the cited paper proves that these transformations
induce the full orientation-preserving orbifold automorphism
action after inner automorphisms have been divided out. A
generating vector is compared modulo $\Aut(G)$ by its rooted,
edge-coloured Cayley graph: breadth-first enumeration with the
ordered generators gives a canonical transition table. Equality
of tables gives the explicit bijection of vertices preserving
every labelled multiplication edge; induction on words proves
that it is the unique target-group isomorphism carrying the
generating vectors to one another. Thus the canonicalization
does not approximate group automorphisms.

To allow conjugacy by negative homeomorphisms, we add the
reflection of the marked sphere. Its effect on the tuple is
\[
 (u_i)_{i=1}^r\longmapsto
 \left(p_{i-1}u_i^{-1}p_{i-1}^{-1}\right)_{i=1}^r,
 \qquad p_{i-1}=u_1\cdots u_{i-1}.
\]
The product of the new tuple is $(u_1\cdots u_r)^{-1}=1$,
as successive cancellation shows. The program computes the
orbits both before and after adding this transformation; it
recovers the published planar counts and finds no fusion.
All nonplanar group/signature pairs in these two genera have
only one positive class in the input classification. Reversing
orientation preserves the group and signature and hence fixes
that class. Consequently the positive counts remain 21 and 55
under conjugacy in the extended group.

For genus two we add all puncture configurations with
$\nu\leq2$. An invariant finite puncture set consists of selected
elliptic orbits and a number $a$ of ordinary orbits. Write
$c_i\in\{0,1\}$ for selection of the $i$th elliptic orbit.
The exact equation is
\[
                  \sum_i c_i\frac{|G|}{m_i}+a|G|=\nu.
                                                               \tag{B.33.1}
\]
The program enumerates every solution. It carries the bits
$c_i$ along under each braid move, exchanging the two bits
when the corresponding two elliptic entries are exchanged.
It then takes orbits of these decorated generating vectors
under the same coloured braids, target automorphisms, and
reflection. These are exactly the decorated topological
classes: an equivariant homeomorphism descends to the quotient
orbifold, and its action on the covering is the compatibility
of its induced automorphism with the epimorphism just described.
Conversely that compatibility lifts the quotient homeomorphism
to an equivariant homeomorphism. Ordinary unlabelled punctures
can be moved to any other such configuration through paths in
the regular part of the quotient; point-pushing the paths is
isotopic to the identity after these additional points are
forgotten and so preserves the original covering kernel.
Thus $a$ records all their remaining data.

There are just two nonplanar genus-two entries. For the trivial
group, there is one configuration for each $\nu$. The other
has group $C_2$ and signature $(1;2,2)$. Its two elliptic
images are the unique involution, and exchanging their small
discs preserves the monodromy after an adjustment of the two
handle paths. The four possible pairs of handle images in
$C_2^2$ are connected by point-pushing either branch point
around a handle: crossing the branch cut adds its nonzero
monodromy to the transverse handle image. These two pushes
change either handle coordinate independently. Hence the
exchange lifts and only the number of selected branch orbits
matters. These elementary moves account for all four vectors
in the supplied data and prove the nonplanar decoration step
used by the code. Equation (B.33.1) and the decorated orbits
give 10 and 23 classes at one and two punctures.

We now derive every section entry from the checked finite
groups. A retained elliptic orbit is the left $G$-set
$G/\langle u_i\rangle$. Its point at the coset
$v\langle u_i\rangle$ has stabilizer
$v\langle u_i\rangle v^{-1}$ and positive generator
$v u_i v^{-1}$. B.3 gives exactly one kernel-conjugacy class
of sections for every point fixed by a specified nontrivial
subgroup $K\leq G$. Such a subgroup is cyclic because it is
contained in the cyclic point stabilizer. Conversely every
such containment gives a section. Thus the exact count is
\[
 \#\Sec(E_K\to K)/\pi
  =\sum_{i:c_i=0}
       \#\{v\langle u_i\rangle:
                    K\leq v\langle u_i\rangle v^{-1}\}.
                                                               \tag{B.33.2}
\]
The formula is used for every subgroup, which the program
enumerates by saturation under adjoining each remaining group
element. It checks that no noncyclic subgroup has a section.
The trivial subgroup has one section, and it is not included
in the nontrivial partial-section totals.

The maximal partial sections are precisely these individual
elliptic points, with $G$ acting by left multiplication on the
cosets. A partial section lies below another exactly when
they have the same point and its subgroup is contained in the
other subgroup. This determines the whole inclusion relation
and proves uniqueness of a maximal extension in the positive
case. Each nontrivial section has centralizer one by B.3.
The total number of nontrivial partial-section classes is
independently checked against
\[
             \sum_{i:c_i=0}\frac{|G|}{m_i}\bigl(\tau(m_i)-1\bigr),
\]
where $\tau(m)$ is the number of positive divisors of $m$.
Finally, for $h=v u_i^b v^{-1}$ the value of $\mu$ at this
section is $b/m_i\pmod\Z$. This lists every power of every
local generator and is exactly the intrinsic multiplier
identified in B.12. Puncture stabilizers are the removed
cyclic groups or the trivial group for ordinary orbits.
All these data are computed with integer permutations and
rational fractions. The output has 109 individual action
rows and 59 independent planar braid-orbit checks.
\end{proof}
